\subsection*{Résultats directeurs et séquence de démonstration}
\label{paper:resultados-rectores}
L’exposé initial anticipe la fonction de la constante de couplage ;
sa construction est effectuée après avoir défini les feuillets d’APP,
l’orientation TRIT, les opérateurs TPK et les prolongements régionaux.
La définition~\ref{lib01:definicion} spécifie le registre ordonné, sa
lecture positionnelle et le système de coordonnées dans lequel ils sont représentés. La
section~\ref{subsec:k-terminal} construit l’extraction terminale et distingue
l’information visible de l’information signée conservée dans l’état.

La fonction arithmétique du registre est établie dans la construction de
\(\alpha\). La section~\ref{subsec:alpha-publicaciones-completas} contient
la représentation analytique complète, sa récurrence à tous les degrés,
la convergence de ses troncatures et l’identification de sa racine à la
précoordonnée engendrée. Sa dépendance à l’égard de cette précoordonnée fait partie de
la définition du lecteur. Cette construction précède l’échelle d’action,
les coordonnées angulaires et les réalisations d’aire et de gravitation.

La fonction de reconstruction s'exprime au moyen de résultats distincts et composables. La proposition~\ref{lib01:marco-ciclico} démontre que l'orbite cyclique du registre constitue une base ; la proposition~\ref{lib01:operadores} détermine quelles réponses permettent de récupérer un opérateur. Le théorème~\ref{lib04:teo-balance} établit la conservation bilatérale pour deux isométries. Le théorème~\ref{lib04:teo-archivo} construit l'archive de profondeur illimitée et son inverse continu. Enfin, les théorèmes~\ref{lib05:teo-biyeccion} et~\ref{lib05:teo-lector} incorporent le raffinement arithmétique et calculent la transformation des observables.

Ces résultats explicitent la hiérarchie de l’article : génération de
l’état, extraction de la constante, réalisations arithmétiques et
incidencielles, conservation des formes et récupération effective de
l’information. Les formules de reconstruction spécifient dans chaque cas le
registre qu’elles reçoivent ; la composition conserve le sens de leurs
domaines, de leurs échelles et de leurs orientations.

% BEGIN CENTRALIDAD_K_20260921
La sélection conjointe de la constante est développée avant sa
reconstruction algébrique : les ordonnancements du répertoire terminal produisent
\(19\,446\) registres ; la compatibilité exceptionnelle en conserve \(79\), et la
compatibilité temporelle détermine un unique registre \(K\). On conserve les
conditions des deux filtres et leur dépendance à l’égard de l’histoire enrichie.
La section~\ref{act21:centralidad-constante} articule cette détermination avec
la clôture arithmétique, les représentations incidencielles et la récupération
opératoire.
% END CENTRALIDAD_K_20260921
